\section{Roadmap Toward Black Hole Stability}
\label{sec:bh_stability}

This section is intentionally \emph{programmatic}. It does not claim a proof of Kerr stability. Instead, it records a possible route by which the elliptic ideas appearing in the Penrose inequality problem could interact with the hyperbolic estimates used in black-hole stability.

\subsection{Target Statement}

The long-term target is the non-linear stability of the sub-extremal Kerr family. Rather than asserting that theorem here, we isolate a collection of analytic subproblems whose resolution would plausibly feed into the Dafermos--Holzegel--Rodnianski style stability framework.

\begin{openproblem}[Hyperbolic realization of the elliptic Penrose structure]\label{op:HyperbolicRealization}
Can one construct a geometrically natural foliation of the black-hole exterior for which the horizon-local elliptic operators governing MOTS/Jang geometry couple coherently to the spacetime wave and Teukolsky equations, with estimates strong enough to produce integrated local energy decay and a nonlinear bootstrap?
\end{openproblem}

\subsection{Three concrete subproblems}

\begin{enumerate}
    \item \textbf{Horizon-local coercivity.} Show that on a preferred family of MOTS sections $\Sigma_t$, the linearized horizon operator admits a coercive estimate of the form
    \[
        \|u\|_{H^1(\Sigma_t)}^2 \le C\Big(\langle \mathcal L_t u,u\rangle + \|u\|_{L^2(\Sigma_t)}^2\Big),
    \]
    uniformly in $t$, after projecting away from the stationary Kerr modes.

    \item \textbf{Propagation from the horizon to the exterior.} Construct a gauge in which the red-shift estimate near the event horizon can be glued to the horizon-local coercive form above and to an $r^p$ hierarchy at infinity. The missing point is not the separate estimates, but a robust \emph{compatibility mechanism} between the horizon geometry and the bulk hyperbolic energy.

    \item \textbf{Low-frequency resolvent control.} Prove that the zero-frequency and near-zero-frequency pieces of the linearized Einstein/Teukolsky system are controlled by the same geometric quantities that govern the horizon-local elliptic operator. This would rule out slowly growing modes that usually dominate the hardest part of the stability argument.
\end{enumerate}

\subsection{A sufficient blueprint}

\begin{proposition}[Conditional blueprint for Kerr stability]\label{prop:ConditionalKerrBlueprint}
Assume the following hold for small perturbations of sub-extremal Kerr initial data:
\begin{enumerate}
    \item[(B1)] a preferred exterior foliation by surfaces $\Sigma_t$ with uniformly controlled geometry;
    \item[(B2)] a coercive horizon-local estimate as in Subproblem~1, modulo the stationary Kerr parameters;
    \item[(B3)] a gluing theorem that combines the red-shift estimate, the horizon-local coercivity, and the usual $r^p$ hierarchy into an integrated local energy decay estimate for the linearized system;
    \item[(B4)] a nonlinear error estimate showing that the quadratic and higher-order terms are integrable in time once the linear estimate from (B3) is available.
\end{enumerate}
Then the standard nonlinear bootstrap for small-data stability of Kerr closes.
\end{proposition}

\begin{proof}
Assumption (B3) supplies the linear integrated decay estimate needed to control the spacetime energy and the radiation field. Assumption (B4) makes the nonlinear source terms perturbative. After projecting out the finite-dimensional Kerr family, the usual continuity argument improves the bootstrap norms, yielding global existence and convergence to a nearby Kerr solution. The proposition is therefore a reduction statement: it isolates where new geometry would have to enter the established hyperbolic framework.
\end{proof}

\subsection{Why this belongs here}

The Penrose program naturally produces horizon-local elliptic structures: MOTS stability operators, Jang-type boundary operators, and weighted inequalities tied to trapped surfaces. The stability problem, by contrast, is controlled by global hyperbolic dispersive estimates. The speculative bridge is that a sufficiently rigid horizon-local coercive theory might control exactly the low-frequency modes that are hardest to exclude in black-hole stability.

\begin{remark}[What is and is not claimed]
Nothing in this manuscript proves black-hole stability, constructs the required foliation, or identifies the linearized Einstein operator with a MOTS stability operator. This section is only a research roadmap, included to document a possible interface between the Penrose inequality ideas developed here and the separate black-hole stability program.
\end{remark}

